\section{Proofs}
\label{app:proofs}

\subsection{Proof of Proposition~\ref{prop:full-spd-geometricization}}

If $u=-Pg$ with $P\succ0$ and $g\neq0$, then
\[
g^\top u=-g^\top Pg<0.
\]
Conversely, suppose $g^\top u<0$ and define $b=-u$.  Then $g^\top b>0$.
Choose an orthonormal basis with $e_1=g/\norm{g}$.  In this basis write
\[
b=\beta e_1+c,\qquad c\perp e_1,
\qquad
\beta=e_1^\top b=\frac{g^\top b}{\norm g}>0.
\]
For $d=1$, set $P=[\beta/\norm g]$.  For $d>1$, define
\[
\widetilde P
=
\begin{pmatrix}
\beta/\norm g & c^\top/\norm g\\
c/\norm g & \gamma I
\end{pmatrix}
\]
in the chosen basis.  Then
\[
\widetilde P(\norm g e_1)=b.
\]
The Schur complement of the upper-left block is
\[
\gamma I-\frac{cc^\top}{\norm g\,\beta}.
\]
Choosing $\gamma>\norm c^2/(\norm g\,\beta)$ makes this complement positive,
so $\widetilde P\succ0$.  Transforming back to the original coordinates gives
an SPD matrix $P$ with $Pg=b=-u$, hence $u=-Pg$.

\subsection{Proof of Corollary~\ref{cor:critical-boundary}}

If $\dd f_\theta=0$, linearity gives $P_\theta\dd f_\theta=0$ for every
positive cometric $P_\theta$.  Thus a pure positive-geometry update is zero.
Any nonzero tangent vector at that point must enter through another module or
through a different modeling convention.

\subsection{Proof of Proposition~\ref{prop:diagonal-expressivity}}

The equation $u=-Pg$ with $P=\diag(p_1,\ldots,p_d)$ and $p_i>0$ is equivalent
coordinate-wise to
\[
u_i=-p_i g_i.
\]
If $g_i=0$, exact expression requires $u_i=0$.  If $g_i\neq0$, exact
expression requires $p_i=-u_i/g_i>0$, equivalently $u_i g_i<0$.

For the residual, the squared Euclidean objective separates:
\[
\inf_{p_i>0}\sum_i (u_i+p_i g_i)^2
=
\sum_i \inf_{p_i>0}(u_i+p_i g_i)^2 .
\]
If $g_i=0$, the term is $u_i^2$.  If $g_i\neq0$ and $u_i g_i<0$, the
positive choice $p_i=-u_i/g_i$ gives zero.  If $g_i\neq0$ and $u_i g_i>0$,
the unconstrained minimizer is negative, so the infimum over $p_i>0$ is the
boundary value $u_i^2$, approached as $p_i\downarrow0$.  If $g_i\neq0$ and
$u_i=0$, the infimum is zero, again approached as $p_i\downarrow0$, but it is
not attained by a strictly positive $p_i$.  This proves the formula and the
attainment statement.

\subsection{Bounded diagonal residual}

\begin{proposition}[Bounded diagonal residual]
\label{prop:bounded-diagonal-residual}
Fix $0<\lambda\le L<\infty$ and let
\[
\Pcal_{\rm diag}^{[\lambda,L]}
=
\{\diag(p_1,\ldots,p_d):\lambda\le p_i\le L\}.
\]
For $g,u\in\R^d$ and $h=I$,
\[
\rho_{{\rm diag}^{[\lambda,L]},I}(u;g)^2
=
\sum_{i=1}^d (u_i+p_i^\star g_i)^2,
\]
where, for $g_i\neq0$,
\[
p_i^\star
=
\min\!\left\{L,\max\!\left\{\lambda,-\frac{u_i}{g_i}\right\}\right\},
\]
and for $g_i=0$ any $p_i^\star\in[\lambda,L]$ gives the same term
$u_i^2$.
\end{proposition}

\begin{proof}
The squared residual separates by coordinate:
\[
\inf_{\lambda\le p_i\le L}\sum_i (u_i+p_i g_i)^2
=
\sum_i \inf_{\lambda\le p_i\le L}(u_i+p_i g_i)^2 .
\]
If $g_i=0$, the coordinate objective is constant and equal to $u_i^2$.  If
$g_i\neq0$, the unconstrained minimizer of the one-dimensional convex
quadratic is $-u_i/g_i$.  Projecting this minimizer onto the interval
$[\lambda,L]$ gives $p_i^\star$, and substituting the coordinate minimizers
gives the stated formula.
\end{proof}

\subsection{Proof of Proposition~\ref{prop:block-expressivity}}

For a block-diagonal $P$, the equation $u=-Pg$ separates over blocks:
\[
u_{B_j}=-P_{B_j}g_{B_j}.
\]
If $g_{B_j}=0$, exact expression requires $u_{B_j}=0$.  If $g_{B_j}\neq0$,
\cref{prop:full-spd-geometricization} applied inside the block gives
existence of $P_{B_j}\succ0$ exactly when
$g_{B_j}^\top u_{B_j}<0$.  Combining the block solutions gives the desired
block-diagonal matrix.
